\begin{Abstract}
    Scientific literature production in higher education institutions shows a significant increase every year. The growing number of publications causes researchers to experience difficulties in finding relevant literature efficiently. Conventional academic search systems that rely on keyword matching are less capable of understanding semantic context and relationships between academic entities. This study aims to design a Knowledge Graph construction pipeline from academic literature corpora, design a Hybrid Vector-Graph Retrieval mechanism within a Knowledge Graph-Based Retrieval Augmented Generation architecture, and evaluate the effectiveness of Knowledge Graph integration in improving the quality of large language model responses. Publication data were obtained from lecturer profiles and publications of the Informatics and Computer Science department at UNESA. The graph construction produced 45,685 nodes and 122,675 relationships with 16 node types covering bibliographic and conceptual entities. Experiments were conducted by comparing three retrieval methods, namely Naive, Subgraph, and Hybrid on 90 evaluation questions divided into factual, relational, and multi-hop categories. Retrieval evaluation used Mean Reciprocal Rank and Hit@K, while answer generation evaluation used the Pairwise LLM as Judge method. The Hybrid method achieved the best performance with an MRR of 0.50 and an Overall Win Rate of 53.9\% against the Naive method. The Subgraph method obtained an MRR of 0.43 with an Overall Win Rate of 45.3\% against Naive. Evaluation on the Traceability and Faithfulness dimensions showed that Knowledge Graph integration helps large language models produce answers that are more traceable to their sources and reduces unverified claims.

    \keywords{Knowledge Graph, Retrieval Augmented Generation, Hybrid Retrieval, Academic Literature Search, Large Language Model}
\end{Abstract}
